\chapter{练习提示}

本章只给出解题路径，不提供可直接抄写的标准答案。一个合格答案通常需要四步：先把问题改写成明确命题，再重构支持理由；随后加入至少一个真正威胁结论的反例，最后说明修改后的结论为何仍值得接受。若题目要求比较理论，还应使用同一案例检验各理论，不能只并排列定义。

\section*{第一章}

\begin{enumerate}
\item 先固定同一个命题、主体和时间，只改变信念为真的方式与主体的理由。“合理但为假”应保留好证据而改变世界；“不合理但为真”应保留结论真值而撤去恰当根据；Gettier 版本则要让主体有看似充分的确证，却由独立偶然因素保证真值。最后列表比较真值、确证、基于关系和运气，避免三个故事同时改变太多变量。
\item 选择Gettier原文中的“工作与十枚硬币”或“福特车与巴塞罗那”案例。证据路径应从主体观察到的材料，经推理到其实际相信的命题；真值路径则说明该命题在世界中为什么为真。两条路径在错误前提或错误对象处发生分离，这个分离而非单纯“用了假命题”才是核心。
\item 对三种修补先各写必要条件，再用相同案例测试。无假前提能处理依赖明显假前提的推理，却可能漏掉无明示推理的环境运气；恰当因果联系能连接事实与信念，却需说明一般命题、否定命题怎样因果支持；反运气条件覆盖更广，但必须交代相关近邻与形成方法怎样固定。结论可以是组合方案，但要说明组合不是临时堵洞。
\item 若坚持史密斯知道，就不能同时保留“知识的成功不能由与主体根据无关的偶然因素造成”这一直觉。他可以弱化知识的反运气要求，或主张确证与真值已经足够，但需承担其理论会把更多幸运真信念算作知识的后果。答案应给出一个新案例，显示这项代价究竟可否接受。
\item 先把现实争论中的核心命题限定到具体时间，分别问：它真吗、主体当时有什么认识资格、推理是否理性、主体是否尽到责任。一个人可以根据良好证据作出合理假判断，也可以因懒惰碰巧说真；制度性误导还会使主体无责却没有知识。至少找出一组评价分离，说明为何不能用“他人很好”替代认识结论。
\item 开头用一句话表明你选择严格定义、解释性模型还是实践概念，并写出选择标准。正文以一个难例展示方案的解释力，再重构最强对手而非弱化版本。结尾应说明方案放弃了什么，例如严格反例免疫、概念统一或日常用法忠实，而不是宣称三项目标都无代价地获得。
\end{enumerate}

\section*{第二章}

\begin{enumerate}
\item 先把怀疑论证写成有效形式，不要把“我无法证明不是缸中脑”直接等同“我不知道有手”。闭合原则负责从知道日常命题推出知道其逻辑后果，排除原则要求知识排除相关错误可能，怀疑前提则否定主体对反怀疑命题的知识。逐项指出拒绝某一前提会失去什么直觉。
\item 摩尔路线保留日常知识并拒绝怀疑前提的优先性；语境主义保留不同对话环境中的知识归属，却需解释标准转换；外在主义允许主体不从内部证明可靠性，却承担“主体无法辨认自己处境”的压力。用同一例子写三种回应，并比较它们回答的是知识是否存在、断言是否合适，还是主体能否反思证明。
\item 停钟案例的真值依赖查询时刻，主要展示介入运气；假谷仓中目标可能真实而环境充满赝品，主要展示环境运气。敏感性问“若命题为假，主体是否仍相信”，安全性问相近情形中以同一方式相信是否容易为假。必须明确可能世界中固定了哪种查看方法与环境，否则判断会随叙述任意变化。
\item 构造案例时先写出方法 (M)、命题 (p) 和相关近邻。要得到敏感但不安全，可让最接近的非-(p) 世界中主体不信 (p)，却存在其他很近的 (p) 或非-(p) 情形使同一方法出错；反向案例则可利用必然或稳定真命题测试敏感性的过强要求。答案重点是解释两种反事实条件为何分开，而非故事离奇。
\item “能力参与”太弱，因为几乎任何成功都包含一些正常感知或推理能力。能力论还需说明真信念在多大程度上因该能力而成功，可用射箭中的阵风类比区分参与与归功。进一步讨论环境支持：若良好设备、同伴校验和制度记录共同成就真信念，归功可以分层，而不必全归个人。
\item 先列出案例中的具体异常，如声音不同步、付款域名变化或路线与路标冲突，再将其与“任何信息都可能被完美伪造”的纯逻辑可能区分。停止核查标准应结合来源独立性、错误成本、时间限制和可逆性，而不是要求绝对确定。最后写出何种新异常会重新开启核查，显示信任是可错且可修正的。
\end{enumerate}

\section*{第三章}

\begin{enumerate}
\item 让主体同时拥有支持结论的好证据与真正驱动其相信的迷信理由，例如医生掌握检测结果却只因占卜接受诊断。命题可能得到确证，主体的信念却没有以该证据为根据，因而信念确证与知识受损。说明反事实或因果测试：撤去迷信理由而保留证据时，主体是否仍会相信？
\item 温和基础主义允许某些经验提供初始但可击败的支持，从而停止回溯；融贯主义把支持放在信念系统的相互关系中。前者需解释初始理由如何与世界相连且不成为武断特权，后者需解释两个同样融贯而只有一个符合世界的系统怎样区分。比较时应让二者面对同一幻觉或孤立群体案例。
\item 恶魔受害者与正常主体拥有相同经验和反思材料，却处在系统误导环境。内在主义倾向判二者同样有主观或责任意义上的确证，可靠主义则在产生知识的客观资格上区分二者。答案须指明分歧落在哪一评价层级，否则会把“受害者无可责”误当作“受害者知道”。
\item 语境主义解释归属者语境改变“知道”的标准；实践侵入主张主体利害可能改变知识或确证本身；纯行动解释则保持知识不变，只提高行动所需复核。新判断题应交叉改变说话者语境、主体利害和行动可逆性，观察哪项变化影响知识归属。若三个理论都能解释原银行案例，区分性变体才是论证中心。
\item 时间线至少包含初始证据、击败信息、复核与最终状态。反驳型击败者支持非-(p)，削弱型击败者攻击来源或方法，误导性击败者虽然本身虚假，也可能在主体尚未解除时损害其理性资格。要分别评价客观知识、主体当时合理性和解除击败后的恢复条件。
\item 可把基于关系写为：主体之所以相信 (p)，必须以支持 (p) 的理由 (R) 为适当的因果、反事实或理性依据。遗忘案例检验过去理由是否仍能维持关系，重复搜索检验熟悉感是否冒充支持，自动推荐检验系统选择是否替代主体评估。若条件排除太多日常记忆知识，应提出保存性机制或程度化版本。
\end{enumerate}

\section*{第四章}

\begin{enumerate}
\item 固定对象与观察条件，再分别改变刺激、经验和分类：正常知觉中对象与经验匹配；错觉有真实对象但某性质呈现失真；幻觉缺少相应对象；误认则可能看清对象却错误分类。逐一定位失败在外界刺激、感知关系、经验内容还是后续判断，避免把所有错误都叫“看错”。
\item 教条主义认为正常经验本身可提供初始理由，无需先证明可靠；析取论强调良好知觉与主观不可分辨的幻觉在认识种类上不同。前者面对幻觉者似乎拥有同样现象根据的压力，后者需解释主体内部无法区分时为何仍可有不同理由。用同一红墙幻觉案例比较二者对主体合理性与知识的评价。
\item 设计同一场景让偏见或训练只改变注意分配，然后检查经验是否真的不同。若主体看到同样特征却下不同结论，变化在判断；若专家主动寻找普通人未查找的线索，变化在搜索策略；若现象呈现本身受概念影响，才涉及经验内容。说明可用何种行为、报告或干预证据区分这些解释。
\item 来源遗忘时，主体也许无法给出当前理由。内在主义会追问记忆内容或背景理由对主体是否可及，外在主义则可诉诸原始获取与保存过程的可靠性。答案还要改变一个条件，例如原始来源后来被揭露不可靠，观察两套评价如何处理保存、击败与责任。
\item “我正在疼”是对体验的认识，“这种感觉属于疼痛”包含概念分类，“疼痛由牙髓炎造成”则是外因解释。前两者也可能受概念能力与注意影响，第三者尤其需要医学证据。用止痛药、幻肢痛或焦虑案例说明内省特权有范围，而不是从体验不可疑直接推出原因判断无误。
\item 先选数学、逻辑或概念关系命题，分别列出学习符号和概念所需的经验、当前支持结论的理由。回应归纳反对时说明推理是否依赖观察多个实例，回应约定反对时说明命题为何不只是语言规则。若最终承认经验参与，也要区分因果上的学习条件与认识上的确证来源。
\end{enumerate}

\section*{第五章}

\begin{enumerate}
\item 还原论要求儿童的证言信任最终由说者可靠性的经验理由支持，却面临儿童尚无足够独立资料的难题；非还原论给予证言初始资格，但需以击败敏感性防止轻信。可用同一个照护者既可靠又偶尔玩笑的案例，说明儿童怎样从默认信任逐步形成来源区分，而非在全信与全疑之间跳跃。
\item 元证据清单不直接判断领域命题，而判断谁更可能有资格。训练说明是否掌握相关方法，记录显示长期表现，透明度允许复核，利益关系提示偏向风险，纠错表现显示面对反证的可靠性。每项都可能被伪装，所以还应写明可验证材料与反例，例如头衔不能替代预测记录。
\item 从最早可定位来源画有向图，页面是节点，转载、改写或独立观察是边。十个页面若都源于同一新闻稿，证据份数接近一而不是十；若来源虽不同却共享同一数据集，独立性仍有限。结论应按命题拆分：多个页面可独立证明传播广泛，却未必独立证明内容真实。
\item 平分论强调同侪身份本身是关于自己可能犯错的新证据；适度降级允许原有一阶理由仍有不对称强度；保持立场则可能诉诸自己掌握而对方未掌握的理由。先确认双方是否真为同侪，再比较信息共享、独立性和过往校准。不可把“我仍确信”当作无需降级的理由。
\item 成员私人相信不等于组织正式接受，正式接受也不自动成为群体知识。选委员会、医院或公司案例，列出授权程序、证据汇集、异议记录、决定发布和后续修正。群体知识至少需正确、适当形成、可由组织调用并对击败者有回应机制；答案可质疑其中某项，但要说明替代条件。
\item 流程要指定谁接收低位者证言、谁复核、谁作决定、谁保存理由，以及申诉和反报复渠道。减少不正义不只是增加发言席位，还要防止证言在记录、议程或权重分配中再次消失。设计审计证据，如异议处理时长、决定修改记录和重复错误率，并讨论透明度与隐私之间的代价。
\end{enumerate}

\section*{第六章}

\begin{enumerate}
\item 一万人中若患病率为 \(1\%\)，先算患者人数与真阳性，再算健康者人数与假阳性，最后用真阳性除以全部阳性。每一步标清分母，就能看到灵敏度或“准确率”不等于阳性后的患病概率。若使用不同数值，必须说明假阳性率而不能只给笼统准确率。
\item 只把先验患病率改为 \(20\%\)，保持检测的灵敏度与特异度不变，重新制作自然频率表。后验上升说明同一似然在不同基础率下提供不同结论。进一步说明参考类为何重要：就诊人群、筛查人群与高风险接触者的基础率不能混用。
\item 让两项检测共享同一批受污染样本、同一设备校准错误或同一训练数据。此时第一次阳性使共同故障更可能，第二次阳性不能按独立证据简单相乘。画出共同原因节点，并说明采用不同实验室、不同原理或盲法复核怎样提高条件独立性。
\item 校准问预测为 \(70\%\) 的事件是否约七成发生，分辨力问模型能否把高低风险拉开，分组表现问这些性质是否在不同群体保持。一个永远报基础率的模型可校准却无分辨力，总体良好也可掩盖小群体系统误差。至少构造一组数值展示三项评价分离。
\item 矩阵的行写行动，列写真实状态，每格列误报、漏报、复核、延迟和不可逆损失。新检测的信息价值是有了结果后采取更好行动所减少的期望损失，减去检测与等待成本。若无论结果怎样行动都不改变，新信息即使提高置信度，决策价值也可能很低。
\item 可选彩票、统计泛化或黑箱模型个案。安全性解释主体在近邻中容易错误，个案联系解释统计支持未连接这个具体对象，行动门槛解释高风险环境需要额外复核。三者可能给出不同结论，因此要说明你是在否定知识、限制断言，还是仅提高行动标准。
\end{enumerate}

\section*{第七章}

\begin{enumerate}
\item 先把真信念固定，避免把知识的价值偷偷归给真理本身。可靠来源可使信念在环境变化中更稳定，使成功归功于能力，使他人更敢传递，也使行动更可问责；这些是不同价值。选择一项作核心并用案例说明，随后检验它是否也适用于无关紧要的真命题。
\item 沼泽化问题问：若真理已经取得，可靠或有根据的过程为何还增值；Gettier问题问：确证真信念为何仍可能不是知识。前者是价值差异，后者是资格缺陷，但二者都迫使理论解释过程与真理的关系。可用同一幸运预测分别问“是否知道”和“若知道又多了什么”。
\item “知道却不理解”可由背诵公式但不能解释或迁移来展示；理想化模型则明知忽略摩擦，却能抓住相关依赖。第二个案例须说明理想化在哪一目的下不扭曲关键结构，以及何种任务会使它失败。这样才能避免把任何有用简化都称为理解。
\item 三项测试可以是反事实回答、跨情境迁移和错误定位，也可加入解释压缩。机械背诵者可能答对原题却在变量变化时失败，模型使用者应能预测干预并指出适用边界。测试是理解的可观察迹象而非定义本身，还需讨论语言表达困难或团队理解的反例。
\item 听者的功劳可能在选择来源、理解与回应击败者，说话者承担生成和表达证据的功劳，制度提供认证、记录与纠错。改变其中一个环节，例如听者随机撞见可靠网页或制度系统性排除异议，观察成就归属怎样变化。结论可接受分布式功劳，而不必把知识全归个人。
\item 先把正确、解释、迁移、校准和修正写成独立可观察任务：例如新题、口头因果图、置信区间和错题修订。为防重复计分，应说明同一表现只进入哪个指标，或怎样设分层权重。再加入一项反游戏设计，防止学生通过套模板获得解释分却没有真正迁移。
\end{enumerate}

\section*{第八章}

\begin{enumerate}
\item 六步必须落实到一条精确命题。保存原回答和时间，打开实际链接而非相信引文外观，核对作者、日期和版本，再逐句判断来源是否支持该命题；独立证据不能只是转载同一材料。最后记录结论是确认、否定还是仍不确定，并写出什么新证据会改变判断。
\item 选择视频、新闻或电商推荐案例，画出内容库、排序信号、界面展示与用户点击的链条。算法可能不修改任何单条真假，却通过曝光频率、默认选项和社会认同线索改变主体可用证据。还要区分排序造成的认识影响与用户原有偏好，提出一个可比较的界面变体。
\item 给两个群体设不同规模，使系统甲乙总体准确率相同，却让一个系统把错误集中在较弱势群体。加入误报和漏报的不同后果后，单一准确率更不能决定公正。答案应同时报告分母、阈值和错误成本，并说明改善一种错误是否会增加另一种。
\item 人工审核者需要看到足够信息、有时间复核、具备相关专业能力，并有真实否决或升级权限。逐项查证，而非把页面上出现“人工参与”当作答案。若审核者只确认格式、受到自动化锚定或否决会受惩罚，该流程就是装饰性监督。
\item 内容溯源回答文件从哪里来，数字签名回答是否由某密钥签发且未被篡改，事实核查评估具体主张，来源可信度则是关于长期表现的可错判断。真实机构也会发布错误，签名也能认证谎言，核查也有范围与时效。用同一伪造或错误文件展示四种工具的互补关系。
\item 先明确“知识归属”要求真信念、可靠能力、理由可及还是责任承担。可靠主义可能偏向稳定运行的系统或团队，能力论追问成功归功于谁，责任主义强调能理解理由和回应批评的主体。用模型更新、人工覆写和机构部署三个环节测试，结论可以是分层归属，但必须说明机器输出与机构知识何时分离。
\end{enumerate}

\section*{第九章}

\begin{enumerate}
\item 先找出待检验的“早已提出”主张和具体文本。语义检验问关键词在原语境是否同义，论证检验问前提与结论结构是否相同，功能检验问概念解决什么问题，历史检验问是否存在传播或影响证据。四项中只要局部相似，就只能得出有限比较，不能推出理论同一。
\item 内部解释要交代文本位置、说话对象、修养或政治语境，不能先用现代定义改写。随后选一个明确现代理论，只比较一项结构，如承认无知的规范作用或认识与行动的联系。结尾列出对象范围、真值角色、主体形成或制度背景上的差异，显示对应关系在哪一步停止。
\item 可选“证成”与“修养”、“知识”与“明”之类不能一一对应的概念。不可通约不是宣布无法对话，而是迫使研究者上移到共同问题，例如判断如何被训练、错误如何纠正。说明保留原概念后，原来的问题分类或评价单位发生了什么改变。
\item 历史同源需要传承、接触或文本影响证据；独立平行需要各自语境中的成形证据并避免虚构传播；规范建构则需公开说明选择与组合的理由。为三类各写一句不同强度的结论，检查是否把“我受二者启发”误写成“二者历史上就是同一理论”。
\item 先用九章各写一个真正不同的问题：知识资格、怀疑与运气、理由与击败、来源、社会依赖、形式更新、认识价值、数字结构和比较方法。然后合并能共享的分析工具，同时保留厚概念、历史角色和责任结构的差异。若九问最后都只得到“多核查来源”，说明分析尚未展开。
\item 理论开头要给出核心条件或解释模型，而不是价值宣言。让西方案例压力测试反例结构，让非西方文本资源挑战问题单位、主体观或认识目的；二者都应可能迫使理论修订。结尾列出明确修正触发器，例如某类反例反复出现、概念无法忠实翻译或制度后果与理论承诺冲突。
\end{enumerate}
